\section{Finite length in characteristic zero}
\label{sec:characteristic-zero}

In this section, we present the first of our two main results, 
which is a method for proving the finite length property assuming that the field has characteristic zero. 
Under this assumption, we will establish the finite length property of the Rado atoms (Example~\ref{ex:rado-graph}) and for the vector atoms (Example~\ref{ex:vector-space-atoms}). 
These are new results. 
We also think that the proof itself, even when applied to get already known results for the equality atoms (Example~\ref{ex:equality-atoms}), is of independent interest and arguably simpler than previously known proofs.
The method that we use will work for all structures satisfying the following condition. 
 
\begin{definition}[Oligomorphic approximation]\label{def:oligomorphic-approximation} 
    We say that a homogeneous structure $\A$ has \emph{oligomorphic approximation} if, 
    for every $d \in \set{1,2,\ldots}$, 
    there exists a family $\Bclass$ of finite substructures of $\A$ such that: 
    \begin{bracketenumerate}
        \item \label{item:oligomorphic-approximation-embedding} 
        every finite substructure $S$ of $\A$ is a substructure of some $\B_S \in \Bclass$; and

        \item \label{item:oligomorphic-approximation-orbits} 
        there is a common (finite) upper bound, for all $\B \in \Bclass$, on the number of orbits in $\B^d$ with respect to automorphisms of $\B$.
    \end{bracketenumerate}
\end{definition}

This is closely related to the notion known as \emph{smooth approximation} in model theory,
developed in~\cite[p.~440]{KLM89} and the subject of the monograph~\cite{CherlinHrushovski_03}.
There we ask for a family $\Bclass$ that does not depend on $d$, 
where in the place of (\ref{item:oligomorphic-approximation-orbits}) we require: 
\begin{description}
    \item[(2')] \label{item:smooth-approximation} 
    \it $\A$ is oligomorphic, and each $\B \in \Bclass$ is homogeneous. 
\end{description}

\begin{theorem}\label{thm:have-oligomorphic-approximation}
    The following structures have oligomorphic approximation: 
    \begin{romanenumerate}
        \item\label{item:eq-has-oligomorphic-approximation} 
        the equality atoms from Example~\ref{ex:equality-atoms};
        
        \item\label{item:vec-has-oligomorphic-approximation}
        the vector atoms from Example~\ref{ex:vector-space-atoms}, for any finite field $\innerfield$;
        
        \item\label{item:smooth-implies-oligomorphic-approximation}   
        any smoothly approximated structure;\footnote{\color{magenta} So that it is homogeneous, we may need to consider the canonical vocabulary.}

        \item\label{item:rado-has-oligomorphic-approximation} 
        the Rado atoms from Example~\ref{ex:rado-graph}.

    \end{romanenumerate}
\end{theorem}

Before proving the above theorem, let us observe that the ordered atoms (Example~\ref{ex:order-atoms}) do not have oligomorphic approximation.

\begin{non-example}\label{non-ex:weak-smooth-approximation}
Consider the rational numbers with the usual order. 
The finite substructures in this case are finite total orders, 
and already for dimension $d=1$, a finite total order of size $n$ has $n$ orbits. 
So no family $\Bclass$ satisfies both conditions (\ref{item:oligomorphic-approximation-embedding}) and (\ref{item:oligomorphic-approximation-orbits}).
\lipicsEnd\end{non-example}

\begin{proof}[Proof of Theorem~\ref{thm:have-oligomorphic-approximation}]
    \begin{romanenumerate}
        \item 
        For the equality atoms, we can simply choose $\Bclass$ to be all finite subsets. 
        These are all homogeneous: we can extend a bijection between subsets to a permutation.

        \item 
        For the vector atoms, we can also choose $\Bclass$ to be all subspaces of finite dimension. 
        These are all homogeneous since we can, by extending a linearly independent set to a basis, extend linear bijections between subspaces to a linear permutation. 
        This argument applies to any finite field $\innerfield$.

        \item
        \textcolor{magenta}{In that case, two tuples in $\B^d$ are in the same orbit if and only if they are in the same orbit of $\A^d$; 
        hence the number of orbits in $\B^d$ is at most that of $\A^d$, which is finite by oligomorphism.}

        \item 
        The most interesting case is the Rado atoms. 
        The witness for oligomorphic approximation %
        % \footnote{But there is no smooth approximation --- see~\cite[Rem.~2.1.2]{CherlinHrushovski_03}. 
        % Essentially, the classification of finite homogeneous graphs~\cite[p.~100]{Gardiner1976} means no such family can satisfy (1).} 
        is a family of \emph{symplectic graphs}: see~\cite[Sec.~8.11]{GR01}.%
        \footnote{We are grateful to Ehud Hrushovski for drawing our attention to this construction.} 
        For every $n \in \set{1,2,\ldots}$ define a finite graph as follows. 
        The set of vertices is the vector space over the two-element field with basis 
        \begin{align*}
            \set{e_1,\ldots,e_n, f_1,\ldots,f_n}.
        \end{align*}
        Since the field has two elements, we can view vertices as subsets of this basis. In this graph, there is an edge between vertices $v$ and $w$ if and only if the sets 
        \begin{align*}
            \setbuild{ i \in \set{1,\ldots,n}}{ $e_i \in v$ and $f_i \in w$ }\\
            \setbuild{ i \in \set{1,\ldots,n}}{ $f_i \in v$ and $e_i \in w$ }
        \end{align*}
        have different sizes modulo two. 
        These graphs satisfy condition~\eqref{item:oligomorphic-approximation-embedding} from Definition~\ref{def:oligomorphic-approximation}, i.e.~every finite graph embeds in some symplectic graph~\cite[Thm.~8.11.2]{GR01}. 
        In Appendix~\ref{sec:appendix-symplectic}, we prove condition~\eqref{item:oligomorphic-approximation-orbits}, i.e.~that the number of orbits of $d$-tuples in symplectic graphs is uniformly bounded by a function of $d$ only. \qedhere
    \end{romanenumerate}
\end{proof}

The main result of this section is the following theorem.

\begin{theorem}\label{thm:weak-smooth-approximation-finite-length}
    If $\A$ has oligomorphic approximation, then it has the finite length property over any field of characteristic $0$.
\end{theorem}

Combining Theorems~\ref{thm:have-oligomorphic-approximation} and~\ref{thm:weak-smooth-approximation-finite-length}, we can get the following results, both old and new, concerning the finite length property.

\begin{corollary}\label{cor:weak-finite-length}
    Over any field of characteristic $0$, the following structures have the finite length property: 
    (\ref{item:eq-has-oligomorphic-approximation}) the equality atoms; 
    (\ref{item:vec-has-oligomorphic-approximation}) the vector atoms; 
    (\ref{item:smooth-implies-oligomorphic-approximation}) all smoothly approximated structures;
    (\ref{item:rado-has-oligomorphic-approximation}) the Rado atoms.
\end{corollary}

As mentioned in the introduction, the finite length property was already known for the equality atoms --- over arbitrary fields even. The results for the vector atoms and the Rado atoms are new. The assumption on characteristic zero is important, at least in the case of the vector atoms where  the finite length property is known to fail over finite fields --- see, most recently,~\cite[Sec.~4.4]{BFKM24}. Later on in this paper, we will prove the result for the Rado atoms again using a different method that works for any field.

The rest of this section is devoted to proving Theorem~\ref{thm:weak-smooth-approximation-finite-length}.

\begin{proof}
    [Proof of Theorem~\ref{thm:weak-smooth-approximation-finite-length}]
Fix a structure $\A$, which has oligomorphic approximation, and  a field of characteristic zero. 
Since the field is fixed, we omit the field subscript and write $\Lin X$ for linear combinations of elements in $X$ that use coefficients from that field. Fix some power $d \in \set{1,2,\ldots}$.  Our goal is to show that 
$
\Lin \A^d
$
 has  the finite length property. For technical reasons, we apply the assumption on oligomorphic approximation not to $d$, but to $2d$, yielding  some class $\Bclass$ of finite structures that satisfies Definition~\ref{def:oligomorphic-approximation}.

We will now proceed in three steps. 
For convenience, we summarise these steps in Figure~\ref{fig:proof-summary-cogless}. 

\begin{figure}
\begin{center}
\begin{tikzcd}[row sep=large]
\text{length of $\Lin \A^d$} 
\ar[d, phantom, "\rotatebox{90}{$\geq$} \text{\quad (Lemma~\ref{lem:reduce-to-finite-substructures})}"']
\\
\displaystyle\sup_{\B \in \Bclass} \text{length of $\Lin \B^d$} 
\ar[d, phantom, "\rotatebox{90}{$\geq$} \text{\quad  (Lemma~\ref{lem:dim-bounds-length})}"']
\\
\displaystyle\sup_{\B \in \Bclass} \text{dimension of } \lineqfun{\Lin \B^d}{\Lin \B^d}
\ar[d, phantom, "\rotatebox{90}{$=$} \text{\quad(Lemma~\ref{lem:bound-on-dim-for-finite-set})}"']
\\
\displaystyle\sup_{\B \in \Bclass} \text{number of orbits in $\B^{2d}$}
\ar[d, phantom, "\rotatebox{90}{$>$} \text{\quad(Definition~\ref{def:oligomorphic-approximation})}"']
\\
\infty.
\end{tikzcd}
\end{center}
\caption{
    \label{fig:proof-summary-cogless}
    Summary of the proof of Theorem~\ref{thm:weak-smooth-approximation-finite-length}.}
\end{figure}

    \begin{lemma}\label{lem:reduce-to-finite-substructures}
        For every $d \in \set{1,2,\ldots}$ we have 
        \begin{align*}
        \text{length of }\Lin \A^d
        \quad \leq \quad  \sup_{\B \in \Bclass} \text{\ length of }\Lin \B^d.
        \end{align*}
   %     where the supremum ranges over finite sets of atoms.
    \end{lemma}
    \begin{proof}
        Consider some chain of equivariant subspaces 
        \begin{align*}
            V_0 \subset V_1 \subset \dots \subset V_n \subseteq \Lin \A^d,
        \end{align*}
        where equivariance is with respect to automorphisms of $\A$. 
        For each $i \in \set{1,\ldots,n}$, choose some vector that is in $V_i$ but not in $V_{i-1}$. 
        Let $S$ be the (finite) set of atoms that appear in these chosen vectors, 
        and choose some $\B \in \Bclass$ which contains all atoms from $S$. 
        Define  
        \begin{align*}
            W_i = V_i \cap \Lin \B^d.
        \end{align*}
        By homogeneity, every automorphism of $\B$ extends to an automorphism of $\A$, 
        and therefore the space $W_i \subseteq \Lin \B^d$ is equivariant with respect to automorphisms of $\B$. 
        The chain of $W_i$'s continues to be strictly growing, 
        since it contains vectors that witness the growth of the original chain. 
        Hence, the new chain forces the length of $\Lin \B^d$ to be at least $n$. 
    \end{proof}


    Thanks to the above lemma, it remains to show that the length of $\Lin \B^d$ is bounded by some number that depends only on $d$. 
    In our proof, this bound will be the number of orbits in $\B^{2d}$. 
    What we have gained by moving from $\A$ to $\B \in \Bclass$ is that our vector spaces now have finite (albeit unbounded) linear dimension, 
    and are acted upon by finite (albeit unbounded) groups. 
    This will let us leverage a basic decomposition result from representation theory called Maschke's Theorem (see e.g.~\cite[Thm.~6.3]{RepTheoryTextbook}).
    In order to be more precise, recall that $\oplus$ stands for the direct sum of vector spaces, i.e.~$V \oplus W$ is the set of pairs $(v,w)$ with $v \in V$ and $w \in W$, where the operations are defined coordinate-wise.

    \begin{fact}[Maschke's Theorem]
        Let $V$ be a finite-dimensional vector space over a field of characteristic zero, equipped with a linear action of a finite group $G$. 
        Then $V$ can be decomposed as 
        \begin{align*}
        V = V_1 \oplus \cdots \oplus V_n,
        \end{align*}
        where each $V_i$ is an equivariant subspace (with respect to the action of $G$) 
        that has length $1$, 
        i.e.~the only equivariant subspaces of $V_i$ are the zero space and the full space $V_i$.
    \end{fact}


    We will use Maschke's Theorem to bound the length of the vector  spaces $\Lin \B^d$ for $\B \in \Bclass$. 
    For two vector spaces $V$ and $W$ equipped with a linear action of the same group $G$, let us write 
    \begin{align*}
        \lineqfun{V}{W}
    \end{align*}
    for the set of all those linear maps from $V$ to $W$ that are \textcolor{magenta}{equivariant} with respect to the action of $G$. 
    (The group and its action are implicit in this notation.) 
    Elements of this set are closed under taking linear combinations, and therefore the set can be seen as a vector space. 
    In particular, it is meaningful to talk about the dimension of this vector space. 
    \begin{lemma}\label{lem:dim-bounds-length}
        Let $V$ be a finite dimensional vector space over a field of characteristic zero, 
        equipped with a linear action of a finite group $G$. 
        Then 
        \begin{align*}
            \text{length of\, $V$} 
            \quad \leq \quad 
            \text{dimension of } \lineqfun{V}{V}.
        \end{align*}
    \end{lemma}
    \begin{proof}
        Apply Maschke's Theorem, yielding a decomposition 
        \begin{align*}
        V = V_1 \oplus \cdots \oplus V_n,
        \end{align*}
        where the subspaces $V_1,\ldots,V_n$ are equivariant and of length $1$, with respect to the action of the group $G$. 
        The length is additive with respect to direct sums (see e.g.~\cite[Prop.~3.17]{RepTheoryTextbook}), i.e.
        \begin{align*}
            \text{length of $V$} 
            = \text{length of $V_1$} + \cdots + \text{length of $V_n$},
        \end{align*}
        so the length of $V$ is equal to $n$. 
        
        We will now consider the right-hand side of the inequality in the statement of the lemma.
        For every $i \in \set{1,\ldots,n}$, we can define an equivariant linear map from $V$ to itself that 
        \begin{itemize}
            \item is the identity on $V_i$, and
            \item maps vectors from the other components to zero.
        \end{itemize}
        This gives us at least $n$ equivariant linear maps from $V$ to itself. 
        None of these maps can be spanned by the others, 
        and hence the dimension is no less than $n$. 
    \end{proof}

    Thanks to Lemmas~\ref{lem:reduce-to-finite-substructures} and~\ref{lem:dim-bounds-length}, 
    the length of the vector space $\Lin \A^d$ is bounded by the dimensions of the vector spaces
    \begin{align}\label{eq:endo-maps-b-d}
        \lineqfun{ \Lin \B^d}{\Lin \B^d}, 
        \tag{$*$}
    \end{align}
    where $\B$ ranges over the family $\Bclass$.
    To complete the proof of the theorem, it still remains to show that these dimensions are bounded by some constant that depends only on $d$. 
    This is done in the following lemma, which completes the proof of Theorem~\ref{thm:weak-smooth-approximation-finite-length}. 

    \begin{lemma}\label{lem:bound-on-dim-for-finite-set} 
        For every $\B \in \Bclass$, 
        the dimension of the vector space in~\eqref{eq:endo-maps-b-d} is 
        precisely the number of orbits in $\B^{2d}$.
    \end{lemma}
    \begin{proof}
        A linear map $\varphi$ in the $\FF$-vector space~\eqref{eq:endo-maps-b-d} amounts to a square matrix indexed by $\B^d$, 
        i.e.~a function
        \begin{align*}
            {\B^d \times \B^d} \to \field ;
            \quad
            (x, y) \mapsto \lambda_y \text{ where } \varphi(1 \cdot x) = \sum_{y'} \lambda_{y'} \cdot y'.
        \end{align*}
        This function must be equivariant with respect to automorphisms of $\B$. 
        This means that inputs in the same orbit must be mapped to the same field element. 
        \textcolor{magenta}{As Reviewer 1 said, note it's $\pi^{-1}(x, y)$.}
        Therefore, to define such a function, we need to choose one element of $\field$ for each orbit of the input. 
        Therefore, the dimension of the space in~\eqref{eq:endo-maps-b-d} is equal to the number of orbits of $\B^{2d}$, 
        under the action of the group of automorphisms of $\B$. 
        \color{black}
    \end{proof}
This dimension is bounded by some constant that depends only on $d$, 
by the assumption of oligomorphic approximation, 
which completes the proof of Theorem~\ref{thm:weak-smooth-approximation-finite-length}.
\end{proof}

The inequalities shown in Figure~\ref{fig:proof-summary-cogless} give us upper bounds on the length. 
In the case of the equality atoms, this bound is exponential in $d$. 
In the case of the vector atoms from Example~\ref{ex:vector-space-atoms}, the bound is the number of linear dependency patterns for $2d$-tuples over a finite field. 
Such a pattern is described by: 
(a)~indicating a subset of the coordinates which is a basis for the tuple; 
and (b)~indicating the basis decompositions for the remaining coordinates. 
This can be done in at most exponentially many ways in $d$, and therefore the overall bound is exponential in $d$.
In Appendix~\ref{sec:appendix-symplectic} we give an upper bound in the case of the Rado atoms, 
also exponential in $d$.
